\section{概念与状态空间}\label{sec:intro}
\src{课件 01 p.13--15、27--34；02 p.4--28；HW1 INTRO1--3}
\par\noindent\begin{tabularx}{\textwidth}{p{.22\textwidth}Y}
\toprule 概念 & 可直接用于短答的解释\\\midrule
理性行动 Rational action & 在已知信息、可用行动与资源条件下，选择预期最能达成目标的行动。依据行动时的信息评价，随机结果不保证每次成功；理性不等于自私或道德好坏。\\
Agent & 接收环境输入，按目标作决策，产生行动或输出的系统；可以是机器人，也可以是软件。\\
四种 AI 定义 & 像人思考、理性思考、像人行动、理性行动。本课偏重理性行动。\\
强/弱 AI & 一般智能的目标 / 解决特定任务的系统。特定任务成功不自动说明具有一般理解。\\
中文屋 & 符号操作产生看似正确的行为，是否等同理解；属于行为与理解关系的哲学讨论。\\
人不擅长的任务 & 例：持续、精确检查百万条记录的重复与规则冲突。说明注意力/记忆负担，以及算法如何用明确规则处理。\\
XOR & 00,11 为 0；01,10 为 1，线性不可分。隐藏层可构造 $h_1=x_1\lor x_2,h_2=x_1\land x_2$，输出 $h_1\land\neg h_2$；纯线性层复合仍线性。\\\bottomrule
\end{tabularx}\par

\ex{搜索建模与计数模板}
\method{题目给现实情境，要求定义搜索或估算规模。}
依次写：状态 $s$、起点 $s_0$、合法动作、后继函数、目标条件、一步代价和路径代价。状态应包含决定未来动作所需信息。
\[
\text{m 个变量，每个 k 值：}\ |S|=k^m;\qquad
\text{固定 b 叉树：第 d 层 }b^d,\quad
N_{0:d}=\frac{b^{d+1}-1}{b-1}.
\]
$b=1$ 时总数 $d+1$；有早停、去重、不等分支时应逐层计数。

\textbf{密码锁例：}3 位、每位 0--9，共 $10^3=1000$ 状态；每步只改变一位 $\pm1$，最多 6 个邻居。每位循环时，从 $x$ 到 $y$ 的最少操作数为
\[
\sum_{i=1}^3\min(|x_i-y_i|,10-|x_i-y_i|).
\]
不循环时为 $\sum_i|x_i-y_i|$。三位十值是三维，不是十维。

\textbf{8-puzzle：}九元组含空格 0；空格和上下左右数字交换，通常每步代价 1。角/边中/中心合法动作数 2/3/4；排列上界 $9!$，固定可解起点可达 $9!/2$。搜索树按历史计节点，原状态图按排列合并。
\checkit{若两个路径到同一棋盘，它们是两个搜索树节点、一个不同状态。BFS 的发现顺序、扩展次数、最短路径、frontier 大小不是同一个量。}
